\documentclass[11pt,reqno,a4paper]{amsart}

\usepackage[T1]{fontenc}
\usepackage{lmodern}
\usepackage{amsmath,amssymb,mathtools}
\usepackage{booktabs}
\usepackage{enumitem}
\usepackage{microtype}
\usepackage[hidelinks]{hyperref}

\newtheorem{theorem}{Theorem}[section]
\newtheorem{proposition}[theorem]{Proposition}
\newtheorem{lemma}[theorem]{Lemma}
\newtheorem{corollary}[theorem]{Corollary}
\theoremstyle{definition}
\newtheorem{definition}[theorem]{Definition}
\theoremstyle{remark}
\newtheorem{remark}[theorem]{Remark}

\newcommand{\comp}[1]{\overline{#1}}
\newcommand{\one}{\mathbf{1}}

\title[The six-color case]{The six-color case of an Erd\H{o}s--Gy\'arf\'as balanced-coloring problem}
\author{Robert Sneiderman}
\date{18 July 2026}
\subjclass[2020]{05C15, 05C35, 05D10}
\keywords{edge coloring, balanced coloring, extremal graph theory, Ramsey theory}

\begin{document}
\hypersetup{
  pdftitle={The six-color case of an Erdos--Gyarfas balanced-coloring problem},
  pdfauthor={Robert Sneiderman}
}

\begin{abstract}
We prove the six-color case of a conjecture of Erd\H{o}s and Gy\'arf\'as:
every six-coloring of the edges of $K_{37}$ has seven vertices whose induced
edges omit at least one color.  For a fixed color, the counterexample
hypothesis gives a graph in which every seven vertices span at most sixteen
edges.  We combine the Kang--Pikhurko extremal theorem with minimum-degree
decompositions to prove clique lemmas on 18, 19, 24, and 25 vertices.  These
lemmas give a 97-edge lower bound for the relevant 31-vertex graphs and then
contradict the edge count of a least color on 37 vertices.  The proof is
non-computational.  This is a fixed-parameter result and does not settle the
conjecture for general $r$.
\end{abstract}

\maketitle

\begin{center}
\fbox{\parbox{0.88\textwidth}{\centering
\textbf{Preprint.}
This manuscript has not completed external mathematical review.}}
\end{center}

\paragraph{AI-use disclosure.}
OpenAI GPT-5.6 Sol was used extensively during proof search, drafting, and
internal checking. OpenAI GPT-5 (Codex) was used for later source review and
release checks. This manuscript has not completed external mathematical
review.

\section{Introduction}

Erd\H{o}s and Gy\'arf\'as asked whether the following assertion holds for
every integer $r\geq3$: whenever the edges of $K_{r^2+1}$ are colored with
$r$ colors, some $r+1$ vertices induce a complete graph on which at least one
color is absent \cite{ErdosGyarfas1999}.  They proved the cases $r=3$ and
$r=4$, observed that the analogous assertion fails for $r=2$, and gave
related constructions on $r^2$ vertices.  The question is also recorded as
Erd\H{o}s Problem 617 \cite{ErdosProblems617}.

We prove the assertion for six colors.

\begin{theorem}\label{thm:main}
For every map $\chi:E(K_{37})\to[6]$, there is a set
$S\subseteq V(K_{37})$ with $|S|=7$ such that
\[
  \chi\!\left(E(K_{37}[S])\right)\ne[6].
\]
\end{theorem}

Theorem~\ref{thm:main} concerns only $r=6$.  No assertion for $r\geq7$ is
proved here.

For one color in a hypothetical counterexample, let $G$ be the graph formed
by the edges of that color.  Every seven vertices span at most sixteen edges
of $G$, since each of the other five colors must occur.  We call this the
local cap.  The argument has three stages.  First, we prove that certain
admissible graphs on 18, 19, 24, and 25 vertices contain $K_6$.  Second, we
peel three such cliques from a 31-vertex graph and reach a triangle-free
13-vertex complement that violates the local cap.  Finally, a least color on
37 vertices reduces by averaging to that 31-vertex obstruction.

The only external graph bounds used below are the Kang--Pikhurko theorem
\cite{KangPikhurko2005} and Brooks's theorem \cite{Brooks1941}.  All finite
case splits in the proof are given explicitly.

\section{Preliminaries}

All graphs are finite and simple.  For a graph $F$, we write $e(F)$ for its
number of edges, $\alpha(F)$ for its independence number, $\tau(F)$ for its
minimum vertex-cover number, and $\delta(F)$ and $\Delta(F)$ for its minimum
and maximum degrees.  The complement of $F$ is $\comp F$, and $F[W]$ denotes
the subgraph induced by $W\subseteq V(F)$.  The Tur\'an graph with $q$ parts
on $n$ vertices is $T_q(n)$, and $t_q(n)=e(T_q(n))$.

\begin{definition}\label{def:admissible}
A graph $F$ is \emph{admissible} if
\[
  e(F[S])\leq16
  \qquad\text{for every }S\in\binom{V(F)}7.
\]
\end{definition}

If $F$ is admissible and $L=\comp F$, then every seven vertices span at
least five edges of $L$.  Admissibility is inherited by induced subgraphs.

We use the following part of the Kang--Pikhurko theorem, together with its
equality characterization \cite[Theorems 1 and 4 and Lemma 5]{KangPikhurko2005}.

\begin{theorem}[Kang--Pikhurko]\label{thm:KP}
Let $q\geq2$, let $n\geq q+3$, and suppose $q\leq(n-1)/2$.  If $L$ is a
$K_{q+1}$-free graph on $n$ vertices that is not $q$-partite, then
\begin{equation}\label{eq:KP}
  e(L)\leq t_q(n)-\left\lfloor\frac nq\right\rfloor+1.
\end{equation}
The equality graphs are those in the construction characterized by Kang and
Pikhurko.
\end{theorem}

We use Brooks's theorem in the following standard form: a connected graph of
maximum degree $D$ has a proper $D$-coloring unless it is complete or an odd
cycle \cite{Brooks1941}.

One fixed equality exclusion will be used to reach the lower endpoint in the
18-vertex lemma.

\begin{lemma}[The 18-vertex endpoint]\label{lem:endpoint18}
Let $R$ be an admissible graph on 18 vertices with $\alpha(R)\leq3$.  If
$\comp R$ is not three-partite, then
\[
  e(R)\geq51.
\]
\end{lemma}

\begin{proof}
The graph $L=\comp R$ is $K_4$-free and not three-partite.  Since
$t_3(18)=108$, Theorem~\ref{thm:KP} first gives
\[
  e(L)\leq108-6+1=103,
  \qquad e(R)\geq50.
\]
Suppose equality holds.  The equality characterization of Kang and Pikhurko
has old part-size sequence $(5,6,6)$ or $(5,5,7)$.  In the complement $R$,
every old part is a clique.  The second sequence therefore contains a
$K_7$, contrary to admissibility.

For $(5,6,6)$, let $N_s$ and $N_t$ be the special and witness parts in
the equality construction, let $x$ be its exceptional vertex, let
$\varnothing\ne A\subsetneq N_s$ be its distinguished subset, and let
$y\in N_t$ be its witness vertex.  In $L$, the vertex $x$ is adjacent to
$A\cup\{y\}$ and to every old part other than $N_s,N_t$, while the edges
between $y$ and $A$ are deleted.  Hence, in $R$, the vertex $x$ is adjacent
to $N_s\setminus A$ and $N_t\setminus\{y\}$, and $y$ is adjacent to $A$.

The parts $N_s,N_t$ have sizes five and six, in either order.  If the
special part has size five,
the witness $K_6$ together with the exceptional vertex spans
\[
  \binom62+5=20
\]
edges of $R$.  If the special part has size six, let $a$, with
$1\leq a\leq5$, be the size of its distinguished subset.  The special
$K_6$, together with either the exceptional vertex or the witness vertex,
spans
\[
  \binom62+\max(a,6-a)\geq18
\]
edges.  Each case exceeds the local cap of sixteen.  Equality is impossible,
so $e(R)\geq51$.
\end{proof}

The next accounting observation will be used repeatedly.

\begin{lemma}[Neighborhood accounting]\label{lem:accounting}
Let $F$ be a graph, let $v$ be a minimum-degree vertex of degree $d$, and put
\[
  N=N_F(v),\qquad U=V(F)\setminus(N\cup\{v\}),
\]
\[
  X=e_F(N,U),\qquad Y=e(F[N]),\qquad A=X+Y.
\]
Then
\begin{equation}\label{eq:accounting}
  X+2Y\geq d(d-1),
  \qquad A\geq\binom d2.
\end{equation}
Consequently,
\begin{equation}\label{eq:nonneighbor-upper}
  e(F[U])\leq e(F)-d-\binom d2.
\end{equation}
If $\alpha(F)\leq s+1$, then $\alpha(F[U])\leq s$.
\end{lemma}

\begin{proof}
Each vertex of $N$ has at least $d-1$ incident edges after its edge to $v$
is removed.  Summing over $N$ counts an $N$--$U$ edge once and an edge of
$F[N]$ twice, which proves the first inequality.  Since
$Y\leq\binom d2$,
\[
  A=(X+2Y)-Y\geq d(d-1)-\binom d2=\binom d2.
\]
The edge decomposition
\[
  e(F)=d+A+e(F[U])
\]
gives \eqref{eq:nonneighbor-upper}.  An independent $(s+1)$-set in $F[U]$
together with $v$ would be an independent $(s+2)$-set in $F$.
\end{proof}

We will use the following elementary averaging fact.

\begin{lemma}[Induced-subgraph averaging]\label{lem:averaging}
Let $F$ have $n$ vertices and $E$ edges, and let $2\leq q\leq n$.  Some
induced $q$-vertex subgraph of $F$ has at most
\[
  E\frac{\binom q2}{\binom n2}
\]
edges.
\end{lemma}

\begin{proof}
Choose a $q$-set uniformly at random.  Each edge of $F$ belongs to the
chosen set with probability $\binom q2/\binom n2$, so the expected number
of induced edges is the displayed quantity.  At least one choice has no
more edges than the expectation.
\end{proof}

We also record the effect of removing a $K_6$.

\begin{lemma}[Clique peeling]\label{lem:peeling}
Let $F$ be admissible, let $B\subseteq V(F)$ span a $K_6$, and suppose
$\alpha(F)\leq s\leq5$.  Then
\[
  \alpha(F-B)\leq s-1,
  \qquad e(F-B)\leq e(F)-15.
\]
\end{lemma}

\begin{proof}
Every vertex outside $B$ has at most one neighbor in $B$, since two such
neighbors together with $B$ would give at least seventeen edges on seven
vertices.  If $I\subseteq V(F)\setminus B$ were an independent $s$-set,
the edges from $I$ to $B$ would meet at most $s<6$ vertices of $B$.  A vertex
of $B$ outside their endpoints would extend $I$ to an independent
$(s+1)$-set.  The edge bound follows by deleting the fifteen edges of $B$.
\end{proof}

The terminal argument below will appear at orders 13 and 14.

\begin{lemma}[Terminal complement]\label{lem:terminal}
Let $L$ be triangle-free with $\alpha(L)\leq5$, and suppose every seven
vertices span at least five edges.  There is no degree-five vertex of $L$
having at least seven nonneighbors.
\end{lemma}

\begin{proof}
Suppose $w$ is such a vertex.  Put $A=N_L(w)$ and let $B$ be its set of
nonneighbors, so $|A|=5$ and $|B|\geq7$.  The set $A$ is independent.  For
$b\in B$, write
\[
  t_b=|N_L(b)\cap A|.
\]
If $b,c\in B$ satisfy $t_b,t_c\leq2$, then the seven-set
$A\cup\{b,c\}$ spans
\[
  t_b+t_c+\one_{bc\in E(L)}
\]
edges.  The local lower bound forces $t_b=t_c=2$ and $bc\in E(L)$.  Thus
the vertices with $t_b\leq2$ form a clique, and triangle-freeness permits at
most two of them.

At least five vertices of $B$ therefore satisfy $t_b\geq3$.  They are
pairwise nonadjacent: if two were adjacent, their neighborhoods in the
independent set $A$ would be disjoint, although their sizes sum to at least
six.  Five such vertices together with $w$ form an independent six-set,
contrary to $\alpha(L)\leq5$.
\end{proof}

\section{The 18-vertex clique lemma}

We now prove the full statement denoted by $P_3$ in the reduction.

\begin{proposition}[$P_3$]\label{prop:P3}
If $R$ is admissible on 18 vertices, $\alpha(R)\leq3$, and $e(R)\leq55$,
then $R$ contains a $K_6$.
\end{proposition}

\begin{proof}
If $\comp R$ is three-partite, one of its three parts has size at least six
and gives a $K_6$ in $R$.  Otherwise Lemma~\ref{lem:endpoint18} gives
$e(R)\geq51$.

Assume that $R$ has no $K_6$.  Choose a minimum-degree vertex $v$, write
$d=d_R(v)$, and let $U$ be its $17-d$ nonneighbors.  Since
$\alpha(R[U])\leq2$, the graph
\[
  L=\comp{R[U]}
\]
is triangle-free.  The absence of a $K_6$ gives $\alpha(L)\leq5$.
Every neighborhood in a triangle-free graph is independent, so
\begin{equation}\label{eq:deltaL5}
  \Delta(L)\leq5.
\end{equation}
Lemma~\ref{lem:accounting} and \eqref{eq:deltaL5} give the following bounds.
\begin{center}
\begin{tabular}{c@{\qquad}c@{\qquad}c@{\qquad}c}
\toprule
$d$ & $|U|$ & lower bound on $e(R[U])$ & lower bound on $e(R)$\\
\midrule
0 & 17 & 94 & 94\\
1 & 16 & 80 & 81\\
2 & 15 & 68 & 71\\
3 & 14 & 56 & 62\\
4 & 13 & 46 & 56\\
\bottomrule
\end{tabular}
\end{center}
Thus $d\geq5$.  Since $e(R)\leq55$, the average degree gives $d\leq6$.

Suppose first that $d=5$.  Now $|U|=12$, and
\eqref{eq:deltaL5} gives $e(R[U])\geq36$.  Write
\[
  e(R)=50+\ell,\qquad
  e_R(N,U)+e(R[N])=10+a,
\]
\[
  e(R[U])=36+k.
\]
Here $1\leq\ell\leq5$, and exact edge accounting gives
\begin{equation}\label{eq:P3-shell}
  a+k=\ell-1,
  \qquad a,k\geq0.
\end{equation}
Let $M=\comp{R[N]}$, put $m=e(M)$, and for $x\in N$ write
$x_x=d_R(x,U)$.  Minimum degree gives
\begin{equation}\label{eq:P3-pointwise}
  x_x\geq d_M(x).
\end{equation}
The sum of the cross-degrees is
\[
  \sum_{x\in N}x_x=(10+a)-(10-m)=a+m.
\]
Summing \eqref{eq:P3-pointwise} gives $a+m\geq2m$, so
\begin{equation}\label{eq:P3-m}
  1\leq m\leq a\leq4.
\end{equation}
The lower bound $m\geq1$ follows because $m=0$ would make $N[v]$ a $K_6$.

The triangle-free graph $L$ is not bipartite, since one part of a
bipartition of its 12 vertices would be an independent six-set in $L$ and
hence a $K_6$ in $R[U]$.  In particular,
\[
  \tau(L)=12-\alpha(L)\geq7.
\]
If $pq\in E(M)$, then
\begin{equation}\label{eq:P3-cover}
  N_R(p,U)\cup N_R(q,U)
\end{equation}
is a vertex cover of $L$.  Indeed, an uncovered edge of $L$ together with
$p,q$ would be an independent four-set in $R$.

There is an edge $pq\in E(M)$ for which $x_p+x_q\leq6$.  If $m\leq2$,
this follows from
\[
  \sum_{x\in N}x_x=a+m\leq6.
\]
If $m=3$, the total is at most seven, and every edge of a three-edge simple
graph has a nonisolated vertex outside its endpoints.  That outside vertex
has positive cross-degree by \eqref{eq:P3-pointwise}.  If $m=4$, the total
is eight, while a four-edge simple graph has at least four nonisolated
vertices.  At least two positive cross-degrees therefore lie outside the
endpoints of any chosen edge.  In every case \eqref{eq:P3-cover} has size at
most six, contrary to $\tau(L)\geq7$.

It remains that $d=6$.  Now $|U|=11$.  Put
\[
  X=e_R(N,U),\qquad Y=e(R[N]).
\]
The local cap on $N[v]$ gives $Y\leq10$, while minimum degree gives
\[
  X+2Y\geq30,
  \qquad X+Y\geq20.
\]
The graph $L=\comp{R[U]}$ is triangle-free and nonbipartite.  Applying
Theorem~\ref{thm:KP} with $q=2$ and $n=11$ gives
\[
  e(L)\leq t_2(11)-\left\lfloor\frac{11}{2}\right\rfloor+1
  =30-5+1=26.
\]
Consequently,
\[
  e(R[U])\geq\binom{11}{2}-26=29
\]
and
\[
  e(R)\geq6+20+29=55.
\]
All inequalities are equalities.  Thus
\[
  e(R)=55,\qquad Y=10,\qquad X=10,\qquad e(L)=26.
\]

Let $M=\comp{R[N]}$.  Then $M$ has five edges.  For each $x\in N$,
minimum degree gives $d_R(x,U)\geq d_M(x)$; equality of the sums gives
\begin{equation}\label{eq:P3-cross-equality}
  d_R(x,U)=d_M(x)
  \qquad(x\in N).
\end{equation}
We also have $\tau(L)=11-\alpha(L)\geq6$.  As above, every edge $pq\in E(M)$
makes the two $U$-neighborhoods a vertex cover of $L$.  Hence
\begin{equation}\label{eq:P3-degree-sum}
  d_M(p)+d_M(q)\geq6
  \qquad(pq\in E(M)).
\end{equation}

We classify a five-edge simple graph $M$ on six vertices satisfying
\eqref{eq:P3-degree-sum}.  If $M$ has a leaf, its neighbor has degree five,
and all five edges form $K_{1,5}$.  If $M$ has no leaf, let $t$ be the
number of nonisolated vertices.  The ten edge-ends give $4\leq t\leq5$.
For $t=4$, the graph is $K_4$ minus one edge, and an edge joining degrees
two and three violates \eqref{eq:P3-degree-sum}.  For $t=5$, every
nonisolated vertex has degree two, so $M=C_5\mathbin{\dot\cup}K_1$, which
also violates \eqref{eq:P3-degree-sum}.  Therefore $M=K_{1,5}$.  Its five
leaves form a $K_5$ in $R[N]$, and adjoining $v$ gives a $K_6$, the final
contradiction.
\end{proof}

\section{The 19-vertex clique lemma}

\begin{proposition}\label{prop:Q19}
If $Q$ is admissible on 19 vertices, $\alpha(Q)\leq3$, and $e(Q)\leq66$,
then $Q$ contains a $K_6$.
\end{proposition}

\begin{proof}
Suppose not.  Choose a minimum-degree vertex $v$, put $d=d_Q(v)$, and let
$U$ be its $18-d$ nonneighbors.  The graph
\[
  L=\comp{Q[U]}
\]
is triangle-free because $\alpha(Q[U])\leq2$.  Since $Q[U]$ has no $K_6$,
$\alpha(L)\leq5$, and hence $\Delta(L)\leq5$.

Write $X=e_Q(N(v),U)$ and $Y=e(Q[N(v)])$.  Lemma~\ref{lem:accounting}
gives
\[
  X+2Y\geq d(d-1),
  \qquad X+Y\geq\binom d2.
\]
For $d=0,1,2,3$, the resulting lower bounds on $e(Q)$ are
\[
  108,\qquad95,\qquad83,\qquad74,
\]
respectively.  The average degree therefore leaves $d\in\{4,5,6\}$.

If $d=4$, then $|U|=14$ and
\[
  e(Q[U])\geq\binom{14}{2}-35=56.
\]
All bounds are equalities.  In particular, $e(Q)=66$, $Y=6$, and $X=0$.
Thus $N[v]$ is an isolated $K_5$, and $L$ is a 5-regular triangle-free
graph on 14 vertices.  Every seven vertices span at least five edges of
$L$.  Any vertex of $L$ has degree five and eight nonneighbors, contrary to
Lemma~\ref{lem:terminal}.

Suppose $d=5$.  Let
\[
  M=\comp{Q[N(v)]},\qquad m=e(M),
\]
and write
\[
  x_a=d_Q(a,U)=d_M(a)+z_a,
  \qquad z_a\geq0,\qquad z=\sum_a z_a.
\]
The 13-vertex graph $L$ satisfies
\[
  e(L)\leq32,
  \qquad \tau(L)=13-\alpha(L)\geq8.
\]
Exact accounting gives
\[
  e(Q)=5+(10-m)+(2m+z)+(78-e(L))
      =93+m+z-e(L)\leq66.
\]
It follows that $m+z\leq5$.  If $m=0$, then $N[v]$ is a $K_6$.  Otherwise
choose $aa'\in E(M)$.  The set
\[
  N_Q(a,U)\cup N_Q(a',U)
\]
is a vertex cover of $L$, for an uncovered edge would extend $a,a'$ to an
independent four-set in $Q$.  Its size is at most
\[
  d_M(a)+d_M(a')+z_a+z_{a'}
  \leq m+1+z\leq6,
\]
contrary to $\tau(L)\geq8$.

It remains that $d=6$.  Now $|U|=12$.  Retain the notation $M,m,x_a,z_a$
and put $Z=\sum_a z_a$.  The local cap on $N[v]$ gives $m\geq5$.  Since
$e(L)\leq30$, exact accounting gives
\[
  e(Q)=6+(15-m)+(2m+Z)+(66-e(L))
      =87+m+Z-e(L)\leq66,
\]
and hence
\begin{equation}\label{eq:Q19-budget}
  Z\leq9-m.
\end{equation}
We also have $\tau(L)\geq12-5=7$.  For every edge $aa'\in E(M)$, the two
$U$-neighborhoods cover $L$, so
\begin{equation}\label{eq:Q19-weight}
  x_a+x_{a'}\geq7.
\end{equation}

We first suppose that $M$ is triangle-free.  If $\Delta(M)=5$, then
$M=K_{1,5}$, and its leaves form a $K_5$ in $Q[N(v)]$ that extends with
$v$ to a $K_6$.  We may therefore assume $\Delta(M)\leq4$.  For every
edge $aa'$ of a triangle-free graph on six vertices,
\[
  d_M(a)+d_M(a')\leq6.
\]
Summing \eqref{eq:Q19-weight} over all edges of $M$ gives
\begin{equation}\label{eq:Q19-sum}
  7m
  \leq\sum_a d_M(a)x_a
  =\sum_a {d_M(a)}^2+\sum_a d_M(a)z_a.
\end{equation}

If $m\geq8$, the first sum is at most $6m$, and the second is at most
$4(9-m)$.  Their sum is smaller than $7m$.  If $m=7$ and
$\Delta(M)\leq3$, the right side is at most $42+6<49$.  If $m=7$ and
$\Delta(M)=4$, the four neighbors of a degree-four vertex are independent,
and the three remaining edges join the sixth vertex to three of them.  The
degree sequence is $(4,3,2,2,2,1)$, so the right side is at most
$38+8<49$.

Let $m=6$.  If $\Delta(M)\leq3$, then
\begin{equation}\label{eq:Q19-square30}
  \sum_a {d_M(a)}^2\leq30.
\end{equation}
To see this, the only degree partitions of sum twelve and maximum part three
with larger square-sum are $(3,3,3,3,0,0)$ and $(3,3,3,2,1,0)$.  The first
would induce $K_4$ on its four nonisolated vertices.  In the second, deleting
the isolated vertex leaves six edges on five vertices; equality in Mantel's
theorem would give degree sequence $(3,3,2,2,2)$, not
$(3,3,3,2,1)$.  Thus the right side of \eqref{eq:Q19-sum} is at most
$30+9<42$.  If $m=6$ and $\Delta(M)=4$, the degree sequence is
$(4,2,2,2,1,1)$.  Equality in \eqref{eq:Q19-sum} would require all three
units of $Z$ to lie on the degree-four vertex.  One of the two edges not
incident with that vertex would then have endpoint weight $2+2<7$.

Let $m=5$.  If $\Delta(M)\leq2$, the right side of
\eqref{eq:Q19-sum} is at most $20+8<35$.  If $\Delta(M)=3$, the sum can
reach 35 only if the degree square-sum is at least 23.  The only possible
partitions are
\[
  (3,3,3,1,0,0),\qquad
  (3,3,2,2,0,0),\qquad
  (3,3,2,1,1,0).
\]
The first is not graphical: three degree-three vertices on four
nonisolated vertices force the fourth degree to be three.  The second is
$K_4$ minus an edge and contains a triangle.  In the third, if the two
degree-three vertices are adjacent, their two remaining neighbor sets among
three vertices intersect and form a triangle.  If they are nonadjacent,
both are adjacent to all three remaining nonisolated vertices, which makes
the two degree-one entries impossible.  Finally, if $\Delta(M)=4$, then
$M$ is $K_{1,4}$ together with an edge from the sixth vertex to one leaf.
That last edge forces all four units of $Z$ onto its endpoints.  A different
center--leaf edge then has endpoint weight $4+1<7$.

This excludes every triangle-free $M$ except the star already handled.
Suppose now that $M$ contains a triangle $T$ but no $K_4$, since a $K_4$ in
$M$ is an independent four-set in $Q$.  Let $W$ be the other three vertices,
and put
\[
  c=e_M(T,W),\qquad q=e(M[W]).
\]
The three vertices of $T$ are pairwise nonadjacent in $Q$.  Their
$U$-neighborhoods cover all twelve vertices of $U$, since an uncovered
vertex would extend $T$ to an independent four-set.  Hence
\begin{equation}\label{eq:Q19-triangle-sum}
  12\leq\sum_{a\in T}x_a
  =6+c+\sum_{a\in T}z_a
  \leq6+c+Z
  \leq12-q.
\end{equation}
Every inequality is an equality.  Thus $q=0$, all slack lies on $T$, and
the three triangle weights sum to twelve.

Since $m=3+c\geq5$, at least one cross-edge exists.  Each vertex of $W$
has at most two neighbors in $T$, because $M$ has no
$K_4$, and it has zero slack.  A triangle vertex incident with a cross-edge
therefore has weight at least five by \eqref{eq:Q19-weight}.  Cross-edges
cannot meet all three vertices of $T$, since their weights would sum to at
least fifteen.  They cannot meet only one, say $a$: every cross-neighbor
would have weight one, forcing $x_a\geq6$, while the opposite triangle edge
would have endpoint-weight sum at most six.

Thus the cross-edges meet exactly two triangle vertices, say $a,b$.  Every
cross-neighbor meets both.  Indeed, if a cross-neighbor met only $a$, then
$x_a\geq6$; since $b$ is also incident with a cross-edge, $x_b\geq5$, while
$x_c\geq d_M(c)=2$, contradicting the triangle weight sum twelve.  Every
cross-neighbor consequently has weight two.  The edge constraints and the
weight sum force
\[
  (x_a,x_b,x_c)=(5,5,2).
\]
The three $U$-neighborhoods partition $U$ into sets $A,B,P$ of sizes
$5,5,2$.  Since $ac\in E(M)$, the set $A\cup P$ covers $L$.  Hence $B$ is
independent in $L$, and $B\cup\{b\}$ is a $K_6$ in $Q$, the final
contradiction.
\end{proof}

\section{The 25- and 24-vertex clique lemmas}

\begin{proposition}\label{prop:H25}
If $H$ is admissible on 25 vertices, $\alpha(H)\leq4$, and $e(H)\leq81$,
then $H$ contains a $K_6$.
\end{proposition}

\begin{proof}
Suppose not, and choose a minimum-degree vertex $v$ of degree $d\leq6$.
If $d\leq5$, its $24-d$ nonneighbors induce a graph with independence
number at most three and, by Lemma~\ref{lem:accounting}, at most
\[
  81-d-\binom d2
\]
edges.  Lemma~\ref{lem:averaging}, with $q=19$, gives the corresponding
expected edge bound below.
\begin{center}
\begin{tabular}{c@{\qquad}c@{\qquad}c}
\toprule
$d$ & number of nonneighbors & expected edge bound\\
\midrule
0 & 24 & $4617/92$\\
1 & 23 & $13680/253$\\
2 & 22 & $4446/77$\\
3 & 21 & $855/14$\\
4 & 20 & $639/10$\\
5 & 19 & $66$\\
\bottomrule
\end{tabular}
\end{center}
Every entry is at most 66, so some induced 19-vertex graph satisfies the
hypotheses of Proposition~\ref{prop:Q19} and contains a $K_6$.

It remains that $d=6$.  Let $N=N_H(v)$, let $R$ be the 18 nonneighbors of
$v$, and put
\[
  m=15-e(H[N]).
\]
The local cap on $N[v]$ gives $m\geq5$.  Each vertex of $N$ needs at least
its missing degree in $H[N]$ as cross-degree into $R$, so at least $2m$
edges join $N$ to $R$.  Therefore
\[
  e(R)\leq81-6-(15-m)-2m=60-m\leq55.
\]
Also $\alpha(R)\leq3$.  Proposition~\ref{prop:P3} supplies a $K_6$.
\end{proof}

The next statement is the input denoted by $P_4$.

\begin{proposition}[$P_4$]\label{prop:P4}
If $P$ is admissible on 24 vertices, $\alpha(P)\leq4$, and $e(P)\leq70$,
then $P$ contains a $K_6$.
\end{proposition}

\begin{proof}
Suppose not.  A minimum-degree vertex $v$ has degree $d\leq5$.  If
$d\leq4$, Lemmas~\ref{lem:accounting} and~\ref{lem:averaging}, with
$q=19$, give the following expected edge bounds among the $23-d$
nonneighbors.
\begin{center}
\begin{tabular}{c@{\qquad}c@{\qquad}c}
\toprule
$d$ & number of nonneighbors & expected edge bound\\
\midrule
0 & 23 & $11970/253$\\
1 & 22 & $3933/77$\\
2 & 21 & $3819/70$\\
3 & 20 & $288/5$\\
4 & 19 & $60$\\
\bottomrule
\end{tabular}
\end{center}
Each is at most 66, so Proposition~\ref{prop:Q19} gives a $K_6$.

Suppose $d=5$.  Write $X=e_P(N(v),U)$ and $Y=e(P[N(v)])$.  Since $N[v]$
does not span a $K_6$, we have $Y\leq9$.  The minimum-degree inequality
$X+2Y\geq20$ therefore gives
\[
  X+Y\geq11.
\]
The 18 nonneighbors of $v$ span at most
\[
  70-5-11=54
\]
edges and have independence number at most three.  Proposition~\ref{prop:P3}
again gives a $K_6$.
\end{proof}

\section{The 31-vertex bound}

\begin{proposition}\label{prop:T31}
If $H$ is admissible on 31 vertices and $\alpha(H)\leq5$, then
\[
  e(H)\geq97.
\]
\end{proposition}

\begin{proof}
Suppose $e(H)\leq96$, choose a minimum-degree vertex $v$ of degree
$d\leq6$, and let $U$ be its $30-d$ nonneighbors.  Then
$\alpha(H[U])\leq4$.

If $d\leq5$, Lemmas~\ref{lem:accounting} and~\ref{lem:averaging}, with
$q=25$, give the corresponding expected edge bound below.
\begin{center}
\begin{tabular}{c@{\qquad}c@{\qquad}c}
\toprule
$d$ & number of nonneighbors & expected edge bound\\
\midrule
0 & 30 & $1920/29$\\
1 & 29 & $14250/203$\\
2 & 28 & $1550/21$\\
3 & 27 & $1000/13$\\
4 & 26 & $1032/13$\\
5 & 25 & $81$\\
\bottomrule
\end{tabular}
\end{center}
Each entry is at most 81.  Proposition~\ref{prop:H25} gives a first
$K_6$, denoted by $B_1$.

If $d=6$, write
\[
  X=e_H(N_H(v),U),\qquad Y=e(H[N_H(v)]).
\]
The local cap on $N_H[v]$ gives $Y\leq10$, while minimum degree gives
$X+2Y\geq30$.  Hence $X+Y\geq20$, and
\[
  e(H[U])\leq96-6-20=70.
\]
Proposition~\ref{prop:P4} gives $B_1$ in this case as well.

Apply Lemma~\ref{lem:peeling} repeatedly.  Removing $B_1$ leaves a
25-vertex graph $Q$ with
\[
  \alpha(Q)\leq4,
  \qquad e(Q)\leq81.
\]
Proposition~\ref{prop:H25} gives a second block $B_2$.  Removing it leaves
a 19-vertex graph $R$ with
\[
  \alpha(R)\leq3,
  \qquad e(R)\leq66.
\]
Proposition~\ref{prop:Q19} gives a third block $B_3$.  Let $C$ be the
remaining 13-vertex graph.  Then
\begin{equation}\label{eq:C13}
  \alpha(C)\leq2,
  \qquad e(C)\leq51.
\end{equation}

The graph $C$ has no $K_6$.  Suppose that $B\subseteq C$ were a $K_6$, and
let $D=C-B$, so $|D|=7$.  If $D$ contains a nonedge $xy$, each of $x,y$
has at most one neighbor in $B$.  Some vertex of $B$ is nonadjacent to both,
forming an independent triple with $x,y$, contrary to \eqref{eq:C13}.  If
$D$ has no nonedge, then $D$ is a $K_7$, contrary to admissibility.

Put $L=\comp C$.  By \eqref{eq:C13}, the graph $L$ is triangle-free.  The
absence of a $K_6$ in $C$ gives $\alpha(L)\leq5$, and hence
$\Delta(L)\leq5$.  Every seven vertices span at least five edges of $L$,
and
\[
  e(L)=\binom{13}{2}-e(C)\geq78-51=27.
\]
If every vertex of $L$ had degree at most four, then $e(L)\leq26$.
Therefore $L$ has a degree-five vertex, which has exactly seven nonneighbors.
This contradicts Lemma~\ref{lem:terminal}.
\end{proof}

\section{Proof of the main theorem}

\begin{proof}[Proof of Theorem~\ref{thm:main}]
Assume for a contradiction that every induced $K_7$ in a six-coloring of
$K_{37}$ uses all six colors.  For each color $i\in[6]$, let $G_i$ be its
color graph.  An independent seven-set in $G_i$ would omit color $i$, so
\begin{equation}\label{eq:color-alpha}
  \alpha(G_i)\leq6.
\end{equation}
On every seven vertices, the other five colors each occupy at least one of
the 21 edges.  Hence
\begin{equation}\label{eq:color-cap}
  e(G_i[S])\leq21-5=16
  \qquad\text{for every }S\in\binom{V(K_{37})}7.
\end{equation}
Thus each $G_i$ is admissible.

Choose a least color graph $G$.  The six color classes partition the
$\binom{37}{2}=666$ edges of $K_{37}$, so
\begin{equation}\label{eq:least111}
  e(G)\leq111.
\end{equation}
We claim that $\delta(G)\leq5$.  If $\delta(G)\geq6$, then
\eqref{eq:least111} forces $G$ to be 6-regular with 111 edges.  No component
is $K_7$, by \eqref{eq:color-cap}, and no component is an odd cycle, since
every component is 6-regular.  Brooks's theorem gives a proper six-coloring
of each component, hence of $G$.  One of the six independent color classes
has at least seven vertices, contrary to \eqref{eq:color-alpha}.

Choose a minimum-degree vertex $v$ of degree $d\leq5$, and let $U$ be its
$36-d$ nonneighbors.  Then $\alpha(G[U])\leq5$.  By
Lemma~\ref{lem:accounting},
\begin{equation}\label{eq:final-U}
  e(G[U])\leq111-d-\binom d2.
\end{equation}
For $d\leq4$, Lemma~\ref{lem:averaging}, with $q=31$, gives the
corresponding expected edge bound below.
\begin{center}
\begin{tabular}{c@{\qquad}c@{\qquad}c}
\toprule
$d$ & number of nonneighbors & expected edge bound\\
\midrule
0 & 36 & $1147/14$\\
1 & 35 & $10230/119$\\
2 & 34 & $16740/187$\\
3 & 33 & $16275/176$\\
4 & 32 & $1515/16$\\
\bottomrule
\end{tabular}
\end{center}
Every value is less than 96.  If $d=5$, then $|U|=31$ and
\eqref{eq:final-U} gives $e(G[U])\leq96$ directly.  In every case, some
induced 31-vertex graph is admissible, has independence number at most five,
and has at most 96 edges.  This contradicts Proposition~\ref{prop:T31}.

The assumed balanced coloring cannot exist.  Therefore some seven vertices
omit a color.
\end{proof}

\enlargethispage{3\baselineskip}
\bibliographystyle{amsplain}
\bibliography{references}

\end{document}
